% LaTeX companion of hw09.typ. Both formats are editable.
\section{Homework 9 --- submitted work}

\chapter{Part A}

The assigned exercises are 5.4: 27, 31; and 5.5: 15, 23, 32(a--d).

\section{5.4 Exercise 27}

The work observes that \(S^{\perp}\) means the least-squares error is orthogonal to \(S\). Thus, for the displayed vector, the least-squares solution is unchanged:

\(x^{=}\begin{pmatrix}
7 \\
11
\end{pmatrix}.\)

\section{5.4 Exercise 31}

For the three points \(\left( {0,3} \right),\left( {1,3} \right),\left( {1,6} \right)\), the line of best fit uses

\(A = \begin{pmatrix}
1 & 0 \\
1 & 1 \\
1 & 1
\end{pmatrix},\quad b = \begin{pmatrix}
3 \\
3 \\
6
\end{pmatrix}.\)

The normal equations recorded in the submission are

\(\begin{pmatrix}
3 & 2 \\
2 & 2
\end{pmatrix}\begin{pmatrix}
c_{0} \\
c_{1}
\end{pmatrix} = \begin{pmatrix}
12 \\
9
\end{pmatrix},\)

and solving them gives \(c_{0} = 3\) and \(c_{1} = \frac{3}{2}\). Hence the submitted line is

\(f(x) = 3 + 3\frac{x}{2}.\)

\section{5.5 Exercise 15}

For the bilinear expression in the exercise, symmetry requires \(b = c\). The submitted positive-definiteness test yields the additional condition

\(d > b^{2}.\)

\section{5.5 Exercise 23}

With

\(\angle\left( {f,g} \right) = \frac{1}{2}\left( {f(0)g(0) + f(1)g(1)} \right),\)

the answer verifies the inner-product properties and gives the orthonormal basis

\(1,\quad 2x - 1.\)

\section{5.5 Exercise 32}

For the weighted integral inner product

\(\angle\left( {f,g} \right) = \frac{1}{2}\int_{- 1}^{1}f(t)g(t)\, dt,\)

the computation in the submitted pages records

when \(n + m\) is even, and

\(\angle\left( {t^{n},t^{m}} \right) = \frac{1}{n + m + 1}\)

when \(n + m\) is odd. Also \(\left\| t^{n} \right\| = \sqrt{\frac{1}{2n + 1}}\). Applying Gram--Schmidt, it writes

\(g_{0} = 1,\quad g_{1} = \sqrt{3}t,\quad g_{2} = \sqrt{\frac{5}{2}}\left( {3t^{2} - 1} \right),\)

\(g_{3} = \frac{1}{\sqrt{\frac{4}{175}}}\left( {t^{3} - \frac{3}{5}t} \right).\)

The associated polynomial sequence is recorded as

\(1,\quad t,\quad\frac{3t^{2} - 1}{2},\quad\frac{5t^{3} - 3t}{2}.\)

\chapter{Part B}

\section{Problem 1}

The submitted solution finds the plane of best fit for the three displayed points by writing its normal-equation system. Its computed coefficient vector is

\(\begin{pmatrix}
\frac{7}{3} \\
1 \\
{- \frac{8}{3}}
\end{pmatrix},\)

which is also plotted on the graph in the original submission.

\section{Problem 2}

\subsection{(a)}

For (i), the proposed expression is not an inner product: the work uses

\(f(x) = x^{2} - 4x + 3,\)

which is nonzero while \(f(1) = f(3) = 0\). For (ii), it verifies the nonnegativity and definiteness of the sum-of-squares evaluation expression, and concludes that it is an inner product.

\subsection{(b)}

For (i), the weighted integral with weight \(x\) fails positive definiteness; the submission gives \(f(x) = \sin(x)\) as the counterexample. For (ii), the weight \(x^{2}\) gives the required symmetric, bilinear, positive-definite inner product.

\section{Problem 3}

Let

\(\angle\left( {f,g} \right) = \int_{- \frac{\pi}{2}}^{\frac{\pi}{2}}f(x)g(x)\sin^{2}(x)\, dx.\)

\subsection{(a)}

The work finds \(\angle\left( {1,x} \right) = 0\) and \(\left\| 1 \right\| = \sqrt{\frac{\pi}{2}}\), then evaluates the remaining displayed polynomial inner products to prepare Gram--Schmidt.

\subsection{(b)}

The submitted orthonormalized functions are recorded approximately as

\(u_{1} = \frac{2}{\sqrt{\pi}},\quad u_{2} = 0.974x,\quad u_{3} = \frac{x^{2} - 4.435}{15.8376}.\)

\subsection{(c)}

For \(f(x) = e^{x}\), the work records

\(\text{proj}_{W{(f)}} = 1.7521u_{1} + 0.2492u_{2} - 8.1697u_{3}\)

and, after expansion,

\(\text{proj}_{W{(f)}} \approx 3.3472 + 0.0236x - 0.5142x^{2}.\)
